\chapter{Общий родитель морщинки и дислокации}
\label{ch:v10-particle-wrinkle-dislocation-reopening}

Длинная цепочка абсолютного масштаба временно приостанавливается, чтобы
вернуться к исходной постановке материи. Ранние образы морщинки ткани и
дислокации вакуума должны быть не метафорами, а двумя следствиями одного
локального родителя.

Дислокационный слой уже строг. Для изометрии
$F:\mathbb C^{90}\to\mathbb C^{105}$ дефектный проектор
\begin{equation}
 P_{\rm cl}=I_{105}-FF^*,
 \qquad \rank P_{\rm cl}=15,
 \qquad \operatorname{ind}F=-15
\end{equation}
даёт минимальный нормированный дефицит
\begin{equation}
 \Delta_{\rm cl}=\frac{15}{105}=\frac17.
\end{equation}
Он топологически защищён, но не задаёт пространственный профиль.

Морщинный слой условно имеет стандартный баланс Деррика
\begin{equation}
 E(L)=AL+\frac BL,
 \qquad L_*=\sqrt{B/A},
 \qquad E_*=2\sqrt{AB}.
\end{equation}
При нормированном свидетеле $A=B=1$ получаются $L_*=1$, $E_*=2$ и
$E''(L_*)=2>0$. Это доказывает совместимость конечного радиуса с
дефектным классом, но не происхождение коэффициентов.

Накопленный родитель остаётся прямой суммой. В координатах радиальной
морщинки и дефектного отклика его локальный гессиан имеет вид
\begin{equation}
 H_{\rm inh}=\begin{pmatrix}2&0\\0&0\end{pmatrix},
 \qquad \rank H_{\rm inh}=1.
\end{equation}
Смешанный блок равен нулю; отсутствуют морфизм локализации пятнадцати
дефектных каналов в ядре профиля и отображение энергии профиля в
спектральный полюс. Условная матрица
$\bigl(\begin{smallmatrix}2&1\\1&1\end{smallmatrix}\bigr)$ положительна и
невырождена, поэтому один настоящий смешанный мост был бы достаточен, но
его коэффициент и оператор пока не унаследованы.

\begin{theorem}[Кинематическая совместимость двух образов]
Топологический дефект индекса $-15$ совместим с конечным стационарным
профилем типа $AL+B/L$. Противоречия между морщинкой и дислокацией нет.
\end{theorem}
\begin{nogocriterion}[Прямая сумма не является частицей]
Совместное наличие индексного класса и отдельного конечного профиля не
доказывает, что дефект локализован этим профилем или что его энергия равна
массе одного спектрального полюса. Для этого необходим ненулевой смешанный
оператор общего родителя.
\end{nogocriterion}
\begin{protocol}\begin{enumerate}
\item совместность требований: \textbf{$8/8$};
\item строгий топологический дефект: \textbf{$7/7$};
\item условный стационарный профиль: \textbf{$3/3$};
\item происхождение общего родителя: \textbf{$0/3$};
\item строгая частица: \textbf{$0/1$};
\item следующий гейт: \textbf{аудит смешанных мостов}.
\end{enumerate}\end{protocol}
Машинный сертификат сохранён в
\path{s2t_v10_particle_wrinkle_dislocation_common_parent_reopening_gate_results.json}.